\documentclass[conference]{IEEEtran}
\IEEEoverridecommandlockouts

\usepackage{amsmath,amssymb,amsfonts}
\usepackage{booktabs}
\usepackage{cite}
\usepackage{graphicx}
\usepackage{url}

\pdfinfo{
  /Title (A Spectrum Sensor for CubeSat Radios)
  /Author (Dylan J. Gormley and Anthony A. Stock)
  /Subject (Post-publication fact-checked revision of the 2021 IEEE CCAAW paper)
  /Keywords (CubeSat; spectrum sensing; spectral correlation analyzer; FPGA)
}
\pdfcatalog{/Lang (en-US)}

\begin{document}

\title{A Spectrum Sensor for CubeSat Radios%
\thanks{NASA's Cognitive Communications Project supported this work. This is a post-publication fact-checked revision dated September 2026, not the IEEE version of record. U.S. Government work not protected by U.S. copyright.}}

\author{
\IEEEauthorblockN{Dylan J. Gormley and Anthony A. Stock}
\IEEEauthorblockA{\textit{Cognitive Signal Processing Branch} \\
\textit{NASA Glenn Research Center}\\
Cleveland, OH, USA \\
dylan.j.gormley@nasa.gov; anthony.stock@nasa.gov}
}

\maketitle

\begin{abstract}
Small-spacecraft missions increasingly use distributed architectures, making awareness of shared radio-frequency spectrum useful for interference mitigation and spectrum sharing. Conventional spectral-correlation analyzers (SCAs), however, can demand more memory and computation than some CubeSat radios can readily provide. This paper develops an FPGA-oriented spectrum sensor based on a low-memory streaming SCA. Under the evaluated signal model, the sensor searches configured symbol-rate and center-frequency candidates for power-of-two-order quadrature amplitude modulation waveforms with square-root-raised-cosine pulse shaping. A power-normalized binary threshold determines whether each candidate symbol rate is present; a center-frequency detector retains the strongest frequency candidate, and FIFO-based scheduling coordinates parallel detector instances. Monte Carlo simulations with two co-channel signals were used to characterize the detector. Across the nine reported scenarios, the symbol-rate detection proportion ranged from 0.89 to 1.0, while the false-alarm proportion per unoccupied symbol-rate candidate ranged from 0.0013 to 0.1213. These metrics assess symbol-rate detection only: the archived scoring code does not test center-frequency correctness. Separate center-frequency samples appeared more dispersed as nominal $E_b/N_0$ decreased. The results support feasibility within the simulated QAM--SRRC conditions, but they do not establish general channel independence, hardware performance, or center-frequency accuracy.
\end{abstract}

\begin{IEEEkeywords}
CubeSat, spectrum sensing, spectral correlation analyzer, FPGA
\end{IEEEkeywords}

\section{Introduction}
By April 2021, the cumulative number of launched nanosatellites and CubeSats had risen rapidly from the levels observed two decades earlier (Fig.~\ref{fig:running_total}). Distributed small-spacecraft concepts spanned different scales: HelioSwarm used a nine-spacecraft baseline~\cite{plice-19}; the four-spacecraft Starling1 technology demonstration was intended to exercise capabilities scalable to at least 100 spacecraft~\cite{sanchez-18}; and separate commercial low-Earth-orbit communications proposals contemplated thousands of satellites~\cite{giambene-18}. The last category comprises large constellations rather than CubeSat formation-flying missions. NASA Ames Research Center also illustrated several notional applications enabled by coordinated small-spacecraft swarms (Figs.~\ref{fig:explorers}--\ref{fig:solarwind})~\cite{cannon-18}.

\begin{figure}[!t]
\centering
\includegraphics[width=\columnwidth]{images/Nanosats_total_2021-04-04.png}
\caption{Cumulative launched nanosatellites and CubeSats through April 4, 2021. Source: Erik Kulu, Nanosats Database~\cite{kulu-21}.}
\label{fig:running_total}
\end{figure}

\begin{figure}[!t]
\centering
\includegraphics[width=\columnwidth]{images/explorers.png}
\caption{Starling-Explorers concept for coordinated flight, communication, and sensors such as ground-penetrating radar or magnetometers; adapted from~\cite{cannon-18}.}
\label{fig:explorers}
\end{figure}

\begin{figure}[!t]
\centering
\includegraphics[width=\columnwidth]{images/dsg.png}
\caption{Starling-DSG concept for visual inspection and monitoring in support of Deep Space Gateway operations; adapted from~\cite{cannon-18}.}
\label{fig:dsg}
\end{figure}

\begin{figure}[!t]
\centering
\includegraphics[width=\columnwidth]{images/lunarnet.png}
\caption{Starling-LunarNet concept for monitoring the lunar surface with an ad hoc mesh network in low lunar orbit; adapted from~\cite{cannon-18}.}
\label{fig:lunarnet}
\end{figure}

\begin{figure}[!t]
\centering
\includegraphics[width=\columnwidth]{images/solarwind.png}
\caption{Starling-SolarWind concept for autonomous radio tomography and in situ measurements from a Sun--Earth L1 halo orbit; adapted from~\cite{cannon-18}.}
\label{fig:solarwind}
\end{figure}

If such distributed architectures become more common, their inter-satellite links may face congestion and interference in shared bands. Spectrum sensing could provide measurements to future dynamic spectrum-management algorithms, but conventional cyclostationary processing can be computationally and memory intensive~\cite{koch-19}. CubeSat radios also operate under restrictive size, mass, power, and compute budgets. This motivates the low-memory streaming spectral correlation analyzer (streaming-SCA) developed in~\cite{gormley-21}.

The streaming-SCA estimates selected points of the spectral correlation function (SCF) for square-root-raised-cosine (SRRC) pulse-shaped QAM waveforms. The evaluated implementation assumes a power-of-two modulation order, a configured search grid, and candidate symbol rates and center frequencies; it therefore is not fully blind. For requested symbol rate $R_s$, let
\begin{equation*}
v=\operatorname{round}\!\left(\frac{R_s N}{f_s}\right)
\end{equation*}
be its zero-based Goertzel/DFT bin, and let $f_{c,k}$ be the $k$th configured physical center-frequency candidate. In zero-based sample notation, the archived estimator is
\begin{equation}\label{eq:sca}
\widehat{S}_{X}^{v}[k]
=\frac{1}{N}\sum_{n=0}^{N-1}
\left|\left(X[n]e^{+j2\pi f_{c,k}n/f_s}\right)\ast h[n]\right|^{2}
e^{-j2\pi vn/N}.
\end{equation}
Here, $n$ is the integer sample index, $N$ is the number of samples used for one estimate, $f_s$ is sample rate, $\ast$ denotes convolution, and $h=[1/2,1/2]$ is the two-tap low-pass filter. The positive channelizer sign matches the archived transmitter's negative carrier convention. The one-based Goertzel call uses bin $v+1$, while $k$ indexes the configured center-frequency list. The quantity $\widehat{S}_{X}^{v}[k]$ is one estimated SCF coefficient~\cite{gormley-21}. Equation~\eqref{eq:sca} corrects the 2021 formula, which omitted frequency normalization and conflated a candidate-list index with physical center frequency.

The spectrum-sensor design searches and structures these SCF estimates. A scheduling algorithm manages requested symbol-rate candidates. Each request is processed by a streaming-SCA instance, then evaluated with the binary test
\begin{equation}\label{eq:sca_test}
\left|\widehat{S}^{v}_{X}[k]\right|^{2}
\underset{H_{0}}{\overset{H_{1}}{\gtrless}}
\beta\left(\frac{1}{N}\sum_{n=0}^{N-1}|X[n]|^{2}\right)^{2}.
\end{equation}
The scalar $\beta$ sets the operating tradeoff between the symbol-rate detection proportion $P_d$ and the false-alarm proportion $P_{fa}$~\cite{gormley-21}. Valid results are stored for a prospective real-time spectrum-management algorithm.

The 2021 study generated a receiver operating characteristic (ROC) curve for the reported two-signal configuration in Table~\ref{tab:classifier}. The repository does not contain the raw ROC observations, and its surviving simulation script uses different signal settings, so the numerical curve cannot be independently reproduced from the archived files.

\begin{table}[!t]
\caption{Reported Two-Signal ROC Configuration}
\label{tab:classifier}
\centering
\footnotesize
\setlength{\tabcolsep}{2.2pt}
\begin{tabular}{@{}cccccccc@{}}
\toprule
Signal & $N$ & $f_s$ & $R_s$ & $f_c$ & $M$ & $r$ & Nominal $E_b/N_0$ \\
 & & (MHz) & (MBd) & (MHz) & & & (dB) \\
\midrule
0 & $2^{12}$ & 1.0 & 0.1 & 0.1 & 32 & 0.2 & 3.0 \\
1 & $2^{12}$ & 1.0 & 0.2 & 0.2 & 32 & 0.2 & 3.0 \\
\bottomrule
\end{tabular}
\end{table}

\begin{figure}[!t]
\centering
\includegraphics[width=\columnwidth]{images/roc.png}
\caption{Archived ROC curve based on 300 trials per value of $\beta$. The red point is approximately $(P_{fa},P_d)=(0.17000,0.64167)$, giving $J=0.47167$. No confidence interval was calculated.}
\label{fig:roc}
\end{figure}

Youden's statistic is a scalar, not an ROC coordinate~\cite{youden-50}:
\begin{equation}\label{eq:J}
J=P_d+(1-P_{fa})-1=P_d-P_{fa}.
\end{equation}
The archived script correctly maximizes $P_d-P_{fa}$ after averaging repeated false-alarm values, but it then uses that averaged-array index directly on the original $\beta$ array. Consequently, the reported $\beta=1.23\times10^{-4}$ is not reliably linked to the red point and is retained below only as the design value reported in 2021, not as a verified optimum. Maximizing $J$ also gives equal weight to sensitivity and specificity; it does not uniquely minimize false alarms.

\section{Digital Design}
This section describes the digital-design additions proposed for autonomous spectrum sensing. The repository contains diagrams and simulation code but no synthesizable HDL, resource report, timing closure, or power measurement; hardware performance is therefore outside the evidence available for this audit.

\subsection{Threshold}
The threshold module receives in-phase and quadrature samples from an analog-to-digital converter. As shown in Fig.~\ref{fig:threshold}, samples are serially magnitude-squared, accumulated, right-shifted, and squared. The reported value of $\beta$ is fixed before synthesis.

\begin{figure}[!t]
\centering
\includegraphics[width=\columnwidth]{images/Threshold.png}
\caption{Block diagram of the threshold module.}
\label{fig:threshold}
\end{figure}

\subsection{Center-Frequency Detector}
The center-frequency detector (CFD) implements the test in \eqref{eq:sca_test}. A streaming-SCA instance computes the test statistic $\ell=|\widehat{S}_{X}^{v}[k]|^2$ in parallel with the threshold module; synchronized outputs are then compared.

The design receives a candidate $v$ and a frequency-control-word (FCW) step size that determines the resolution of searched $k$ values. For each $v$, the CFD evaluates all candidate center frequencies, selects the largest statistic, and emits its $k$ only if that maximum exceeds the threshold. This process produces a detected $(v,k)$ pair, as illustrated in Fig.~\ref{fig:cfd}.

\begin{figure}[!t]
\centering
\includegraphics[width=\columnwidth]{images/CFD.png}
\caption{Block diagram of the center-frequency detector module.}
\label{fig:cfd}
\end{figure}

\subsection{Spectrum Sensor}
The spectrum-sensor module stores requested $v$ values in a FIFO and assigns them to available CFD instances. A process polls the instances to determine which are ready for input. The proposed scheduler intends to present at most one result to the demultiplexer per cycle; this requires deterministic equal CFD latency or output arbitration, neither of which can be verified without the HDL. The demultiplexer writes resulting $(v,k,\ell)$ tuples to a result FIFO for a prospective spectrum-management process. Fig.~\ref{fig:sca} illustrates the design.

\begin{figure}[!t]
\centering
\includegraphics[width=\columnwidth]{images/SCA.png}
\caption{Block diagram of the spectrum-sensor module.}
\label{fig:sca}
\end{figure}

\section{Results}
The archived study averaged 300 trials for each scenario in Table~\ref{tab:performance}. The table preserves the published point estimates, but their labels require two qualifications. First, the scoring code tests symbol-rate membership only; it does not test whether the selected center frequency is correct. Thus, $P_d$ is the detected fraction of true symbol-rate candidates and $P_{fa}$ is the false-positive fraction of unoccupied symbol-rate candidates, rather than a whole-system probability. Second, the transmitter peak-normalizes each waveform and calls the noise routine without measuring its average input power; it also adds noise independently to each signal before summing them. The listed $E_b/N_0$ values are therefore nominal routine inputs, not verified received ratios. New quantitative conclusions require a corrected, seeded rerun.

\subsection{Characterization}
\begin{table}[!t]
\caption{Reported Symbol-Rate Detection Performance}
\label{tab:performance}
\centering
\small
\setlength{\tabcolsep}{5pt}
\begin{tabular}{@{}ccccc@{}}
\toprule
Experiment & Nominal $E_b/N_0$ (dB) & $r$ & $P_d$ & $P_{fa}$ \\
\midrule
Baseline & 15.0 & 0.35 & 1.0    & 0.0013 \\
1        & 12.0 & 0.35 & 1.0    & 0.0284 \\
2        & 9.0  & 0.35 & 1.0    & 0.0584 \\
3        & 6.0  & 0.35 & 1.0    & 0.0854 \\
4        & 3.0  & 0.35 & 0.9967 & 0.1146 \\
5        & 0.0  & 0.35 & 0.89   & 0.1213 \\
6        & 15.0 & 0.30 & 1.0    & 0.0108 \\
7        & 15.0 & 0.25 & 1.0    & 0.0108 \\
8        & 15.0 & 0.20 & 1.0    & 0.0150 \\
\bottomrule
\end{tabular}
\end{table}

At the reported design value $\beta=1.23\times10^{-4}$, the symbol-rate detection proportion was 0.89 in the lowest nominal $E_b/N_0$ scenario and at least 0.9967 in the other eight scenarios. The reported false-alarm proportion ranged from 0.0013 to 0.1213. Because the threshold-to-ROC mapping is not recoverable and no confidence intervals or independent validation set were retained, these values neither verify an optimal threshold nor confirm ROC accuracy.

The archived scoring code explicitly counts a result as detected when its symbol-rate candidate is correct, even if its center-frequency candidate is inaccurate. The five retained center-frequency data files contain 300 samples each and show visibly greater dispersion in the more impaired cases, but the repository does not preserve a mapping from every file to its complete simulation configuration. The evidence therefore supports only the qualitative statement that center-frequency dispersion appeared to increase as nominal $E_b/N_0$ decreased in these simulations; it does not establish center-frequency accuracy or a rolloff trend.

\begin{figure}[!t]
\centering
\includegraphics[width=\columnwidth]{images/EbN0_fc_histo.png}
\caption{Archived center-frequency histogram from five 300-sample data files. The original ordinate is labeled ``Probability Density,'' but the plotting code and bin-width normalization were not retained, so its scale cannot be verified. Dispersion appears greater as nominal $E_b/N_0$ decreases.}
\label{fig:histo}
\end{figure}

\subsection{Timing}
To illustrate the proposed scheduling scheme, consider two CFDs and input/output FIFOs of depth two. In Fig.~\ref{fig:timing_in}, the upstream system supplies $v$ candidates. The scheduler assigns $A_1$ and $A_2$ to the two CFDs. Candidates $A_3$ and $A_4$ wait until a CFD is available; $A_5$ is not accepted because the input FIFO is full, as indicated by \texttt{inFifo\_full}. After processing, the results for $A_1$ and $A_2$ enter the result FIFO, allowing $A_3$ and $A_4$ to be assigned.

\begin{figure}[!t]
\centering
\includegraphics[width=\columnwidth]{images/SCA_input.png}
\caption{Timing diagram for pushing $v$ candidates to the input FIFO.}
\label{fig:timing_in}
\end{figure}

In the output example (Fig.~\ref{fig:timing_out}), the downstream module initially makes no pop request. When the results for $A_3$ and $A_4$ arrive, the oldest results in the full result FIFO are discarded. Two subsequent pop requests drive the remaining results sequentially to their outputs.

\begin{figure}[!t]
\centering
\includegraphics[width=\columnwidth]{images/SCA_output.png}
\caption{Output timing example. When the result FIFO is full, older results are discarded in favor of newer values.}
\label{fig:timing_out}
\end{figure}

\section{Conclusion}
This work proposed an FPGA-oriented spectrum-sensor architecture built around the streaming-SCA, with candidate scheduling, threshold detection, and result buffering. The archived Monte Carlo outputs characterize symbol-rate detections for simulated co-channel QAM--SRRC signals, and the timing diagrams describe the proposed FIFO scheduling behavior. They do not validate center-frequency accuracy, a fully blind search, implementation resource use, timing performance, or operation over a general RF channel.

Future work should first restore an executable and versioned simulation, normalize average signal power, add channel noise once after combining signals, define a candidate grid independently of ground truth, retain random seeds and raw observations, and rerun the ROC analysis with a correct mapping from each averaged operating point to $\beta$. Symbol-rate and center-frequency errors should then be scored separately with uncertainty intervals. Only after that recalibration should the detector be assessed on additional signal classes. Hardware claims would require synthesizable HDL, resource and timing reports, and power measurements.

\section*{Acknowledgment}
The authors thank Erica K. Henrichsen for the graphic designs illustrated in this paper.

\newpage
\IEEEtriggeratref{5}
\bibliographystyle{main/IEEEtran}
\bibliography{main/IEEEabrv,main/ccaaw_2021}

\end{document}
